\documentclass[11pt]{article}
\usepackage{amsmath,amssymb}
\title{A uniform gauge quotient for mixed harmonic shears}\author{}\date{}
\begin{document}\maketitle
\paragraph{Scope.}
Retain the original complete four-dimensional cube $[0,R]^4$, all its
boundary links and plaquettes, and the relative interior gauge group.
For each link direction $\mu$, let $a_\mu(x)$ be independent of $x_\mu$
and discrete harmonic in the three other coordinates, with its actual
boundary values. Work first in one fixed Cartan color. The sum of these
four shears is an actual Maxwell-harmonic, divergence-free boundary
extension. No boundary values are set to zero. Write
\[
 E(a)=\sum_p|(da)_p|^2,\qquad
 E_{\rm sep}(a)=\sum_{\mu\ne\nu}\sum_{p\parallel\mu\nu}
                         |\delta_\nu a_\mu|^2 .
\]
The second sum counts the original component derivatives, so that
$E_{\rm sep}(a)$ is the sum of the four individual shear energies.
In general $E_{\rm sep}(a)$ is NOT controlled by $E(a)$: the field
$a_1=x_0,a_0=x_1$ is pure gauge.

\paragraph{1. A uniform quotient, with the actual boundary gauge retained.}
There is a multiaffine scalar polynomial $q$ on the whole cube such that
\[
 \boxed{E_{\rm sep}(a-dq)\le C E(a),}                         \tag{1}
\]
with $C$ independent of $R$ and of the original boundary values.
The polynomial has degree at most one in each of the four coordinates;
its gradient space has dimension fifteen, including constant vectors.
Subtracting $dq$ preserves independence in each component's own
coordinate and transverse discrete harmonicity. It is an honest
complete boundary gauge tangent, not a deletion of physical boundary data.

Here is a proof that keeps the lattice boundary.
On the unit continuous cube define the constant-coefficient operator
\[
 {\cal A}u=(\{\partial_\mu u_\nu-\partial_\nu u_\mu\}_{\mu<\nu},
                  \partial_0u_0,\ldots,\partial_3u_3).
\]
It is complex elliptic: if $\xi\wedge v=0$ for $\xi\in\mathbb C^4
\setminus\{0\}$, then $v=c\xi$; the remaining symbol equations
$c\xi_i^2=0$ imply $c=0$.
Its kernel consists exactly of gradients of multiaffine polynomials.
Indeed curl zero gives $u=\nabla q$ on the cube, and
$\partial_i u_i=0$ gives $\partial_i^2q=0$ for every $i$.

Theorem 1.3 of Diening--Gmeineder, \emph{Sharp Trace and Korn Inequalities
for Differential Operators}, Potential Analysis 63 (2025), 169--222,
DOI 10.1007/s11118-024-10165-1, gives
$\|u\|_{W^{1,2}}\le C(\|{\cal A}u\|_2+\|u\|_2)$ on this cube.
Removing the finite-dimensional kernel gives
\[
 \inf_{v\in\ker{\cal A}}\|\nabla(u-v)\|_2
                         \le C\|{\cal A}u\|_2.                \tag{2}
\]
For completeness, take the $L^2$-orthogonal complement of the kernel.
If its full $W^{1,2}$ estimate failed, normalize a countersequence to
unit $W^{1,2}$ norm. Rellich compactness gives an $L^2$ limit in the
kernel and its orthogonal complement, hence zero. The preceding
theorem then forces the $W^{1,2}$ norm to zero, a contradiction.
This argument imposes no boundary condition on $u$.

Interpolate each lattice component multilinearly in its three transverse
coordinates and constantly in its own coordinate. This is the cubical
Whitney one-form interpolation, which here is also a continuous
$W^{1,2}$ vector field. Its own-coordinate derivatives vanish exactly.
On each elementary cube, its curl is the bilinear interpolation of
the original plaquette curls in the two remaining coordinates.
The tensor-product mass matrices for linear interpolation have fixed
positive eigenvalues; consequently
\[
 \|\nabla I(a-dq)\|_{L^2([0,R]^4)}^2\asymp E_{\rm sep}(a-dq),
 \quad
 \|\operatorname{curl}Ia\|_{L^2([0,R]^4)}^2\asymp E(a).
\]
Boundary plaquettes have multiplicity between one and four, so these
comparisons include them with uniform constants. Multiaffine gradients
interpolate exactly. Rescale (2) from the unit cube to $[0,R]^4$;
both derivative norms have the same scaling. This proves (1).

\paragraph{2. Complete scalar connections and cross terms.}
Set $\widetilde a=a-dq$. The original charged vector Hessian plus
gauge-fixing operator has on component $\mu$ the raw forward phase
\[
 \theta^\mu_\nu(x)=
 a_\nu(x)+a_\nu(x+\hat\mu)+a_\mu(x+\hat\nu)-a_\mu(x).           \tag{3}
\]
This is the sum of the four literal single-shear phases of the harmonic shear leading kinetic bound calculation;
it is a first-background-derivative identity, not a finite nonabelian
superposition. The scalar gradient gauge $a_\mu$ removes the last
difference, leaving $a_\nu(x)+a_\nu(x+\hat\mu)$. Ghost phases are $2a_\nu$.
The exact original insertion check
the finite algebraic calculation verifies (3),
including mixed polarizations and a multiaffine pure-gauge addition,
on $R=2,3$ with zero residual.

For each component grid, a closed magnetic current is the sum of the
four polarization currents. Gauge phases cancel on the closed path,
and
\[
 \left|\sum_{\nu=0}^3J_\nu\right|^2\le4\sum_\nu|J_\nu|^2.
\]
Thus the actual killed/reflected path estimates apply separately
to $\widetilde a_\nu$ and then sum. When the current coordinate is the
Neumann coordinate, transverse harmonicity supplies exactly the
stopped-martingale cancellation, INCLUDING the original killing
boundary values and endpoint term. All other current coordinates
are Dirichlet. We obtain for each scalar determinant
\[
 0\le(\log\det L)''[\widetilde a,\widetilde a]
       \le C\log(2+R)E_{\rm sep}(\widetilde a)
       \le C\log(2+R)E(a).                                   \tag{4}
\]
This controls mixed scalar terms by positivity of the path-current
quadratic form. It does not infer cross bounds for the indefinite
complete determinant from its separate signed restrictions.

The spin insertion is a bounded local linear function of the TOTAL
plaquette flux. The original product Green bound
$|G_\mu(x,y)|\le C(1+|x-y|)^{-2}$ therefore gives
$|\operatorname{Tr}(GS'GS')|\le C\log(2+R)E(a)$ directly.
The diagonal cosine and neutral partners cost $CE(a)$.
Mixed scalar/spin traces and the trace of $S''$ vanish by component
structure. All these are the full finite-boundary partners of the harmonic shear leading kinetic bound calculation.

\paragraph{3. Complete leading response, and exact limits.}
The complete boundary-source Ward identity annihilates $dq$ in both
slots of the leading Hessian. Its flat first derivative vanishes, so
there is no omitted boundary-path acceleration. Consequently (4) and
the retained spin, neutral and ghost multiplicities prove
\[
 \boxed{|F_{\rm quad}''[a,a]|
       =|h_a^*\widehat K h_a|
       \le C\log(2+R)E(a).}                                  \tag{5}
\]
The equality is the uncut harmonic matching of the complete harmonic kinetic matching calculation.
In canonical source units the right side is multiplied by $B^{-1}$.
Flat global color invariance extends (5) to three color-valued shears:
the color Hessian is diagonal and identical in the three colors.
This uses only a second derivative at the flat background, not a
commuting finite-source family.

The gauge quotient is proved on the actual complete cube and does
not alter any retained reference boundary constraint except by its
honest simultaneous gauge transformation. Raw determinant constants
and the source-zero scalar are not removed.
This extends the growing family of the harmonic shear leading kinetic bound calculation to all four shear
polarizations. General harmonic one-forms whose components depend
on their own coordinates are NOT covered. Neither cut self-energy
insertions, the complete first nonlinear coefficient $C_R''$, nor the
original nonlinear remainder or S94 are proved here.
The next indispensable estimate is a controlled decomposition of
general retained harmonic boundary data, or a direct current estimate
for that data, with its original boundary response retained.

\paragraph{A precise obstruction to extending this by gauge choice alone.}
For every sum of shears,
$\delta_\mu\delta_\nu(da)_{\mu\nu}=0$: in the first term
$\delta_\mu a_\nu$ the $\nu$ derivative vanishes, and conversely.
This identity is gauge invariant, and fails for genuine harmonic data.

Write $\Delta_i f(x)=f(x+\hat i)+f(x-\hat i)-2f(x)$.
On the integer lattice put
\[
 f=x_0^3x_1^3,\quad
 p_2=x_2^2/2,\quad p_4=(x_2^4-x_2^2)/24,\quad
 \psi=f-p_2(\Delta_0+\Delta_1)f
                  +p_4(\Delta_0+\Delta_1)^2f.
\]
Since $\Delta_2p_2=1$, $\Delta_2p_4=p_2$, and
$(\Delta_0+\Delta_1)^3f=0$, one has $\Delta\psi=0$ exactly.
Set $a_0=\delta_1^-\psi$, $a_1=-\delta_0^-\psi$, $a_2=a_3=0$.
Backward divergence vanishes by commutation, and $\Delta a=0$.
Thus restriction to a complete cube is an actual Maxwell-harmonic
boundary extension (modulo its original interior gauge).
Its curvature is $F_{01}=-(\Delta_0+\Delta_1)\psi$, and
\[
 \delta_0\delta_1F_{01}(1,1,1,x_3)=-48.
\]
A cube large enough to retain this stencil therefore cannot gauge this
datum into the shear sum. This is not an obstruction to a full kinetic
bound; it shows exactly why an additional harmonic-current estimate,
rather than gauge alignment alone, is needed beyond (5).
\end{document}
